% ================================================================
% ECOVISION AI | PRESENTACION CON RESULTADOS DEFINITIVOS (5 DIAPOSITIVAS)
% Compilacion: PDFLaTeX (Overleaf o TeX Live)
% Resultados: MODO_RAPIDO=False | test independiente de 379 imagenes
% ================================================================
\documentclass[aspectratio=169,10pt]{beamer}
\usepackage[utf8]{inputenc}
\usepackage[T1]{fontenc}
\usepackage[spanish,es-nodecimaldot]{babel}
\usepackage{lmodern}
\usepackage{graphicx}
\usepackage{booktabs}
\usepackage{array}
\usepackage{tikz}
\usetikzlibrary{positioning,arrows.meta,shapes.geometric}

% PALETA
\definecolor{Azul}{HTML}{11253B}
\definecolor{Turquesa}{HTML}{008F85}
\definecolor{VerdeClaro}{HTML}{EAF7F4}
\definecolor{GrisClaro}{HTML}{F2F5F8}
\definecolor{Texto}{HTML}{273746}
\definecolor{Secundario}{HTML}{647485}
\definecolor{Naranja}{HTML}{F3B75E}

\setbeamercolor{background canvas}{bg=white}
\setbeamercolor{normal text}{fg=Texto}
\setbeamercolor{frametitle}{fg=Azul}
\setbeamercolor{structure}{fg=Turquesa}
\setbeamercolor{block title}{bg=Azul,fg=white}
\setbeamercolor{block body}{bg=GrisClaro,fg=Texto}
\setbeamercolor{footline}{bg=Azul,fg=white}
\setbeamerfont{frametitle}{series=\bfseries,size=\Large}
\setbeamertemplate{navigation symbols}{}
\setbeamertemplate{itemize item}{\color{Turquesa}\small$\blacktriangleright$}
\setbeamersize{text margin left=0.70cm,text margin right=0.70cm}
\setbeamertemplate{footline}{%
  \leavevmode\hbox{%
  \begin{beamercolorbox}[wd=.84\paperwidth,ht=2.7ex,dp=1.1ex,leftskip=.7cm]{footline}%
    \scriptsize\textbf{EcoVision AI} \hspace{.8em} | \hspace{.8em} Deep Learning para Vision
  \end{beamercolorbox}%
  \begin{beamercolorbox}[wd=.16\paperwidth,ht=2.7ex,dp=1.1ex,rightskip=.7cm plus1fil]{footline}%
    \hfill\scriptsize\insertframenumber/5
  \end{beamercolorbox}}%
}

% ======== CAMPOS EDITABLES ========
\newcommand{\Integrantes}{David Santiago Cifuentes Grimaldo\\[3pt]Juan David Jim\'enez}
\newcommand{\AccCNN}{56,20\%}
\newcommand{\AccResNet}{\textbf{78,63\%}}
\newcommand{\FunoCNN}{53,63\%}
\newcommand{\FunoResNet}{\textbf{76,16\%}}
\newcommand{\MejorModelo}{ResNet18 (F1 validaci\'on: 73,99\%)}
\newcommand{\HallazgoPrincipal}{81 errores en 379 im\'agenes; principales confusiones entre vidrio y pl\'astico.}

% Esta macro muestra la figura si ya esta disponible, y un recuadro si falta.
% Copiar luego las graficas desde el directorio artifacts/ del notebook.
\newcommand{\figurapendiente}[3]{%
  \IfFileExists{#1}{%
    \includegraphics[width=#2,height=#3,keepaspectratio]{#1}%
  }{%
    \fcolorbox{Turquesa}{GrisClaro}{%
      \parbox[c][#3][c]{\dimexpr#2-2\fboxsep-2\fboxrule\relax}{%
        \centering\color{Secundario}\small\textbf{Figura por insertar}\\[4pt]
        \texttt{\detokenize{#1}}%
      }%
    }%
  }%
}

\begin{document}

% ------------------ DIAPOSITIVA 1 ------------------
\begingroup
\setbeamercolor{background canvas}{bg=Azul}
\begin{frame}[plain]
  \vspace{-0.35cm}
  \begin{tikzpicture}[remember picture,overlay]
    \fill[Turquesa] ([xshift=9.7cm]current page.north west) rectangle (current page.south east);
    \fill[Naranja] ([xshift=.73cm,yshift=-1.22cm]current page.north west)
      rectangle ([xshift=2.45cm,yshift=-1.30cm]current page.north west);
    \draw[white,opacity=.24,line width=1.2pt]
      ([xshift=10.9cm,yshift=-2.1cm]current page.north west) circle (1.25cm);
    \draw[white,opacity=.24,line width=1.2pt]
      ([xshift=10.9cm,yshift=-2.1cm]current page.north west) circle (.78cm);
  \end{tikzpicture}
  \color{white}
  \vspace{.9cm}
  {\small\bfseries PROYECTO FINAL \;\textbar\; DEEP LEARNING PARA VISI\'ON}\par
  \vspace{1cm}
  {\Huge\bfseries EcoVision AI}\par
  \vspace{.2cm}
  {\Large Clasificaci\'on inteligente de residuos}\par
  \vspace{.45cm}
  {\normalsize Comparaci\'on de una CNN propia y ResNet18\par
   para el reconocimiento visual de seis tipos de residuos.}
  \vfill
  {\small\Integrantes}\par
  \vspace{.5cm}
\end{frame}
\endgroup

% ------------------ DIAPOSITIVA 2 ------------------
\begin{frame}{El problema y los datos}
  \vspace{.05cm}
  \begin{columns}[T,onlytextwidth]
    \begin{column}{.51\textwidth}
      {\color{Turquesa}\bfseries ¿QUE BUSCAMOS RESOLVER?}\par\vspace{.14cm}
      Clasificar autom\'aticamente una fotograf\'ia de un residuo para
      apoyar su identificaci\'on y separaci\'on.
      \vspace{.35cm}
      \begin{block}{Dataset p\'ublico: Garbage Classification}
        \begin{tabular}{@{}lr@{\hspace{.36cm}}lr@{}}
          Cart\'on & 403 & Vidrio & 501 \\
          Metal  & 409 & Papel  & 594 \\
          Pl\'astico & 480 & Basura & 137 \\
        \end{tabular}
        \par\vspace{.13cm}
        \textbf{Total: 2.524 im\'agenes v\'alidas.}\\
        {\scriptsize Se descartaron 3 copias exactas.}
      \end{block}
    \end{column}
    \begin{column}{.45\textwidth}
      {\color{Turquesa}\bfseries PROCESAMIENTO}\par\vspace{.17cm}
      \begin{block}{1. Control de calidad}
        Validaci\'on de archivos y eliminaci\'on de duplicados exactos.
      \end{block}
      \vspace{.03cm}
      \begin{block}{2. Divisi\'on estratificada}
        1.766 para entrenamiento, 379 validaci\'on y 379 prueba\\{\scriptsize(70\% / 15\% / 15\%, estratificado).}
      \end{block}
      \vspace{.03cm}
      \begin{block}{3. Aumento de datos}
        Rotaci\'on, cambios de color y volteo solo en entrenamiento.
      \end{block}
    \end{column}
  \end{columns}
  \vspace{.2cm}
  {\scriptsize\color{Secundario}\textbf{Observaci\'on:} el conjunto est\'a desbalanceado; ``basura general'' tiene menos ejemplos que las otras clases.}
\end{frame}

% ------------------ DIAPOSITIVA 3 ------------------
\begin{frame}{Arquitecturas y metodolog\'ia}
  \vspace{-.05cm}
  \begin{columns}[T,onlytextwidth]
    \begin{column}{.47\textwidth}
      \begin{block}{MODELO A \;\textbar\; CNN desde cero}
        \begin{itemize}
          \item Tres bloques convolucionales.
          \item 30.630 par\'ametros; 12 \`epocas.
          \item Entrenada sin pesos preentrenados.
        \end{itemize}
      \end{block}
    \end{column}
    \begin{column}{.47\textwidth}
      \begin{block}{MODELO B \;\textbar\; ResNet18}
        \begin{itemize}
          \item Red residual con pesos de ImageNet.
          \item 3.078 entrenables; 8 \`epocas.
          \item Clasificador de seis categor\'ias.
        \end{itemize}
      \end{block}
    \end{column}
  \end{columns}
  \vspace{.17cm}
  \begin{beamercolorbox}[rounded=true,sep=.15cm,wd=\textwidth]{block body}
    \centering\small \textbf{Configuraci\'on:} PyTorch + AdamW + CrossEntropy ponderada.
    Selecci\'on mediante \textbf{F1 macro de validaci\'on}, sin utilizar test.
  \end{beamercolorbox}
  \vspace{.22cm}
  \centering
  \figurapendiente{figuras/f1_validacion.png}{.98\textwidth}{2.52cm}
\end{frame}

% ------------------ DIAPOSITIVA 4 ------------------
\begin{frame}{Resultados y an\'alisis de errores}
  \vspace{-.06cm}
  \begin{columns}[T,onlytextwidth]
    \begin{column}{.43\textwidth}
      {\color{Turquesa}\bfseries EVALUACI\'ON EN TEST}\par\vspace{.16cm}
      \renewcommand{\arraystretch}{1.37}
      \begin{tabular}{@{}lcc@{}}
        \toprule
        \textbf{M\'etrica} & \textbf{CNN} & \textbf{ResNet18} \\
        \midrule
        Accuracy & \AccCNN & \AccResNet \\
        F1 macro & \FunoCNN & \FunoResNet \\
        \bottomrule
      \end{tabular}\par
      \vspace{.30cm}
      {\small\textbf{Modelo elegido:}}\par\vspace{.08cm}
      {\small\MejorModelo}\par
      \vspace{.22cm}
      {\small\color{Secundario} \HallazgoPrincipal}
      \vspace{.14cm}
      {\scriptsize\color{Secundario} Cart\'on: F1=91,53\%. Basura general: F1=56,41\%.}
    \end{column}
    \begin{column}{.54\textwidth}
      \centering
      \figurapendiente{figuras/matriz_confusion.png}{\linewidth}{4.62cm}\\[.13cm]
      {\scriptsize\color{Secundario}Matriz de confusi\'on del modelo ganador.}
    \end{column}
  \end{columns}
  \vspace{.14cm}
  {\scriptsize\color{Secundario}Evaluaci\'on final en test (379 im\'agenes). ResNet18 acierta 298 casos.}
\end{frame}

% ------------------ DIAPOSITIVA 5 ------------------
\begin{frame}{Aplicaci\'on, limitaciones y conclusiones}
  \vspace{-.1cm}
  {\color{Turquesa}\bfseries DE LA RED NEURONAL A UNA HERRAMIENTA UTILIZABLE}\par
  \vspace{.13cm}
  \centering
  \begin{tikzpicture}[node distance=0.6cm, >=Latex]
    \node[draw=Turquesa,fill=VerdeClaro,rounded corners=3pt,minimum width=2.85cm,minimum height=.88cm,align=center] (a) {\textbf{Fotograf\'ia}\\\scriptsize Usuario carga imagen};
    \node[draw=Turquesa,fill=VerdeClaro,rounded corners=3pt,minimum width=2.85cm,minimum height=.88cm,align=center,right=of a] (b) {\textbf{Modelo elegido}\\\scriptsize Preprocesamiento e inferencia};
    \node[draw=Turquesa,fill=VerdeClaro,rounded corners=3pt,minimum width=2.85cm,minimum height=.88cm,align=center,right=of b] (c) {\textbf{Resultado}\\\scriptsize Clase y confianza};
    \draw[->,thick,Turquesa] (a)--(b);
    \draw[->,thick,Turquesa] (b)--(c);
  \end{tikzpicture}
  \vspace{.25cm}
  \begin{columns}[T,onlytextwidth]
    \begin{column}{.54\textwidth}
      \begin{block}{Aplicaci\'on Streamlit: pr\'oximo paso}
        El modelo se export\'o a \texttt{best\_model.pth}. Falta verificar
        la aplicaci\'on con nuevas fotos y desplegarla.
      \end{block}
      \vspace{.08cm}
      \begin{block}{Conclusi\'on principal}
        ResNet18 mejor\'o el F1 macro en \textbf{22,53 puntos}
        frente a la CNN; persisten errores de separaci\'on de materiales.
      \end{block}
      \vspace{.02cm}
      {\scriptsize\color{Secundario}\textbf{Limitaciones:} una partici\'on aleatoria y datos de estudio; no se ha validado con residuos mezclados o im\'agenes externas.}
    \end{column}
    \begin{column}{.42\textwidth}
      \centering
      \figurapendiente{figuras/gradcam.png}{\linewidth}{2.97cm}\\[.12cm]
      {\scriptsize\color{Secundario}Grad-CAM: regiones asociadas con la predicci\'on.}
    \end{column}
  \end{columns}
\end{frame}

\end{document}
